\documentclass[leqno]{article}
% Shared mark scheme preamble for resources.danieljsmith.org: the pack
% preamble, unchanged, then the scheme's own layer. A scheme's .tex is
% rendered from its YAML (scripts/README.md, Mark schemes), never edited.
\input{\CurrentFilePath/pack_preamble.tex}

% The page: equal side margins, so the table runs the full width of the
% header rule. The footers' right edge follows the text's.
\geometry{left=0.8in, right=0.8in}
\makeatletter
\renewcommand{\djs@footR}[1]{%
  \put(\dimexpr1in+\hoffset+\oddsidemargin+\textwidth\relax,
    -\dimexpr\paperheight-0.32in\relax){\makebox[0pt][r]{#1}}}
\makeatother

\usepackage[export]{adjustbox}
\usepackage{tabularray}
\definecolor{msGrid}{RGB}{176, 172, 164}
\definecolor{msTint}{RGB}{245, 243, 239}

% The header on every page, the front page included, and the booklet label
% under the contents-page title. \djsQuestionPageStyle stays empty.
\newcommand{\djsSchemeHeader}[1]{%
  \djsPackHeader[markscheme]{\textit{Mark scheme} - #1}%
  \renewcommand{\djsPackTitle}{#1}%
  \renewcommand{\djsBookletLabel}{Mark scheme}%
}

% The front page: the contents heading, spaced as tocloft spaces the packs'
% "Questions", then the contents table. A question's name and page number
% both link to its first page.
\newcommand{\msContentsHeading}{%
  \vskip\cftbeforetoctitleskip
  {\noindent\cfttoctitlefont\contentsname\par}%
  \vskip\cftaftertoctitleskip}
\NewTblrEnviron{mscontents}
\SetTblrInner[mscontents]{
  colspec = {Q[wd=26mm, l] X[l] Q[wd=16mm, c] Q[wd=14mm, r]},
  width = \linewidth,
  rowsep = 4pt, colsep = 6pt,
  hlines = {0.4pt, msGrid}, vlines = {0.4pt, msGrid},
  hline{2} = {2.2pt, scarlet},
  row{1} = {font=\small\bfseries},
  row{Z} = {font=\bfseries},
  column{2} = {font=\small},
}
\newcommand{\msContentsQuestion}[1]{\hyperref[q:#1]{Question #1}}
\newcommand{\msContentsPage}[1]{\hyperref[q:#1]{\pageref*{q:#1}}}
\newcommand{\msSep}{\hspace{0.9em}}

% The conventions pages' layout: headings in the documents' own bold, tables
% ruled and padded as the scheme's, roomy lists.
\newcommand{\msHead}[1]{%
  \par\vspace{18pt}%
  {\noindent\bfseries #1\par}%
  \vspace{6pt}}
\NewTblrEnviron{mstblr}
\SetTblrInner[mstblr]{rowsep = 6pt, colsep = 6pt,
  hlines = {0.4pt, msGrid}, vlines = {0.4pt, msGrid}}
\newenvironment{msList}
  {\begin{itemize}[leftmargin=1.2em, itemsep=5pt, topsep=3pt, parsep=0pt]}
  {\end{itemize}}

% A question's first page: the link target and the PDF bookmark, as
% \questionitem gives the packs.
\newcommand{\msQnum}{}
\newcommand{\msQuestion}[1]{%
  \clearpage
  \renewcommand{\msQnum}{#1}%
  \phantomsection\label{q:#1}%
  \addcontentsline{toc}{subsection}{\protect\djsTocEntry{#1}}}

% The scheme's table: Question 16mm, Scheme the rest, Marks 16mm, Notes
% 44mm; every cell ruled in msGrid; the scarlet rule under the head row. A
% long table has no caption or "continued" foot; a continuation page opens
% with a muted line above the repeated head row. The rules across a part's
% rows, the alternative's tint and the total row are the rendered table's.
\NewTblrTheme{ms}{%
  \DefTblrTemplate{firsthead}{default}{}%
  \DefTblrTemplate{middlehead,lasthead}{default}{%
    {\small\itshape\color{djsMuted}Question \msQnum\ continued\par}%
    \vspace{6pt}}%
  \DefTblrTemplate{firstfoot,middlefoot,lastfoot}{default}{}%
}
\NewTblrEnviron{msscheme}
\SetTblrOuter[msscheme]{long, theme=ms, entry=none, label=none}
\SetTblrInner[msscheme]{
  colspec = {Q[wd=16mm, l] X[l] Q[wd=16mm, c] Q[wd=44mm, l]},
  width = \linewidth,
  rowsep = 6pt, colsep = 6pt,
  rowhead = 1,
  vlines = {0.4pt, msGrid},
  hline{1} = {0.4pt, msGrid}, hline{2} = {2.2pt, scarlet},
  hline{Z} = {0.4pt, msGrid},
  column{1} = {font=\bfseries}, column{3} = {font=\bfseries},
  column{4} = {font=\small},
  row{1} = {rowsep=4pt, font=\small\bfseries},
}

% Two marks on a row: each note starts with its code.
\newcommand{\msNote}[2]{\textbf{#1}\enspace #2}
% A full-row figure: the worked solution's own picture, scaled uniformly and
% only down, to the table less the Question column. The picture is drawn
% with the line width it has in the solutions document, at its list depth
% there (the optional argument, 1 by default: a part's working), so one
% sized by \linewidth keeps its size; tabularray sets the cell's \linewidth
% to the cell's measured width, which would shrink it at every pass. The
% widths are measured once, in a solution box at the packs' text width
% (pack_preamble.tex's margins, 0.8in and 1in).
\newlength{\msSpanW}
\makeatletter
\newcommand{\ms@keepLW}[1]{%
  \expandafter\xdef\csname ms@solnLW@#1\endcsname{\the\linewidth}}
\AtBeginDocument{%
  \setlength{\msSpanW}{\dimexpr\textwidth-16mm-28pt\relax}%
  \setbox\z@\vbox{%
    \begin{minipage}{\dimexpr\paperwidth-1.8in\relax}
      \begin{djsSolution}
        \ms@keepLW{0}%
        \begin{enumerate}\item \ms@keepLW{1}%
        \begin{enumerate}\item \ms@keepLW{2}%
        \begin{enumerate}\item \ms@keepLW{3}%
        \end{enumerate}\end{enumerate}\end{enumerate}
      \end{djsSolution}
    \end{minipage}}}
\newcommand{\msFigure}[2][1]{%
  \adjustbox{max width=\msSpanW, max totalheight=0.6\textheight, valign=t}{%
    \hsize=\csname ms@solnLW@#1\endcsname
    \linewidth=\hsize
    #2}}
\makeatother
% A figure mark's feature, in the Scheme column.
\newcommand{\msFeature}[1]{\textit{#1}}
% An alternative's title row.
\newcommand{\msAlternative}[1]{\textit{Alternative:} #1}
% Drafts only: a queried mark's code, the query's why under its note, and a
% flag in its part's place.
\newcommand{\msQueried}[2]{#1~\(\to\) #2?}
\newcommand{\msQuery}[1]{{\itshape\color{djsMuted}Query: #1}}
\newcommand{\msFlag}[2]{\textbf{#1:}\enspace #2}


\djsSchemeHeader{Proof by Induction - Divisibility}
\djsSitePath{Core Pure}{Proof by Induction}

\djsLastUpdated{25 October 2026}
\begin{document}
\djsSetQuestionCount{8}
\djsAddMarks{1}{6}
\djsAddMarks{2}{6}
\djsAddMarks{3}{6}
\djsAddMarks{4}{6}
\djsAddMarks{5}{6}
\djsAddMarks{6}{6}
\djsAddMarks{7}{6}
\djsAddMarks{8}{6}
\vspace*{6pt}
\djsContentsTitle
\vspace{30pt}
\msContentsHeading
\begin{mscontents}{}
Question & Marks by part & Marks & Page \\
\msContentsQuestion{1} & & 6 & \msContentsPage{1} \\
\msContentsQuestion{2} & & 6 & \msContentsPage{2} \\
\msContentsQuestion{3} & & 6 & \msContentsPage{3} \\
\msContentsQuestion{4} & & 6 & \msContentsPage{4} \\
\msContentsQuestion{5} & & 6 & \msContentsPage{5} \\
\msContentsQuestion{6} & & 6 & \msContentsPage{6} \\
\msContentsQuestion{7} & & 6 & \msContentsPage{7} \\
\msContentsQuestion{8} & & 6 & \msContentsPage{8} \\
\SetCell[c=2]{r} Total & & 48 & \\
\end{mscontents}
\newpage
% The conventions pages, shared by every pack's mark scheme: one source of
% truth, \input after the front page, with no strand or pack logic. The keys
% on the first side, the rules on the second. The scheme layer defines
% \msHead, msList and mstblr.
\vspace*{6pt}%
{\noindent\cfttoctitlefont Reading the mark scheme\par}%
\vspace{10pt}%
These rules apply throughout, unless a note says otherwise.

\msHead{Marks}
Every mark is worth one, so \textbf{M1 A1} on a row is two marks. Only marks
coded ft follow through.\par
\vspace{10pt}%
\begin{mstblr}{
  colspec = {Q[wd=15mm, l] X[l]},
  width = \linewidth,
  hline{2} = {2.2pt, scarlet},
  column{1} = {font=\bfseries},
  row{1} = {font=\small\bfseries, rowsep=4pt},
}
Code & Meaning \\
M1 & Method: a correct method, applied to this question. A slip doesn't cost it if the method is clear. A formula must be used, not just quoted. \\
dM1 & Dependent method: as M1, but only if the method mark before it was earned, unless its note names another. \\
A1 & Accuracy: a correct answer or step. It needs the method mark it follows. \\
A1ft & As A1, following through from an earlier answer that the note names. \\
A1~AG & The last mark of a part whose answer is printed in the question: complete working, with no errors. \\
B1 & A correct value, fact or statement. It needs no method mark. \\
B1ft & As B1, following through from an earlier answer that the note names. \\
E1 & A written explanation, interpretation or conclusion. Concluding a proof by induction needs every earlier mark. \\
\end{mstblr}

\msHead{Abbreviations}
ft and AG qualify a code in the Marks column. The others appear in the
notes.\par
\vspace{10pt}%
\begin{mstblr}{
  colspec = {Q[wd=17mm, l] X[l] Q[wd=15mm, l] X[l]},
  width = \linewidth,
  column{1} = {font=\bfseries}, column{3} = {font=\bfseries},
}
ft & follow through & AG & answer given in the question \\
cao & correct answer only, with no follow-through & cso & correct solution only: no errors in the part \\
oe & or an equivalent form & awrt & any answer that rounds to this value \\
soi & seen, or implied by later working & isw & ignore working after a correct answer \\
SC & special case: credit outside the main method & their & the student's own earlier value \\
condone & a slip that doesn't cost the mark & [\;\,] & working the mark doesn't need \\
s.f. & significant figures & d.p. & decimal places \\
\end{mstblr}

\newpage
\msHead{Reading the scheme}
\begin{msList}
\item Each row is a line of working, with the marks it earns beside it.
  Where a row carries two marks, each note starts with its code.
\item \textbf{[\textit{n}]} ends each part with its total.
\item A shaded \textit{Alternative} is a second method for the same total.
  Mark an answer by the method it follows.
\item Any other valid method earns the equivalent marks, unless the question
  names a method or says \textit{hence}.
\item A sketch earns one B1 for each feature named below its figure.
\end{msList}

\msHead{Accuracy and answers}
\begin{msList}
\item Unless a question says otherwise, give answers to 3 significant
  figures, and angles in degrees to 1 decimal place. Angles in radians take
  3 significant figures.
\item A final answer with more figures is accepted if it rounds to the
  scheme's value; one with fewer is not. An answer spoiled by rounding too
  early loses its accuracy mark.
\item Where a question asks for an exact answer, a decimal loses the final
  accuracy mark. Otherwise either is accepted.
\item Wrong or missing units cost nothing unless the question asks for them.
\item A correct answer with no working earns full marks, unless the question
  needs the working: show that, prove, hence, a method it names, or an exact
  answer. A wrong answer with no working earns nothing. Calculator answers
  are treated the same way.
\item A correct answer, once reached, isn't spoiled by later working (isw).
\end{msList}

\msHead{Value of $\boldsymbol{g}$}
\begin{msList}
\item Take $g = 9.8\,\mathrm{m\,s^{-2}}$ unless a question states a value,
  which is then the only one accepted. Answers that use $g = 9.8$ may be
  given to 2 or 3 significant figures.
\item Where a question doesn't state $g$, an answer worked with $g = 10$
  throughout earns the same marks, keeping any value printed in the
  question. The worked solutions list the final answers with $g = 10$.
\end{msList}


\msQuestion{1}
\begin{msscheme}{
  hline{3-7} = {2-Z}{0.4pt, msGrid},
  hline{8} = {0.4pt, msGrid},
  row{8} = {font=\bfseries},
}
Question & Scheme & Marks & Notes \\
1 & {\(n=1:\ 1^5-1=0\), divisible by 5} & B1 & {Needs the value 0 and a statement that it is divisible by 5, or true for \(n=1\).} \\
& {Assume true for \(n=k\): \(k^5-k=5m\)\\[4pt] Considers \((k+1)^5-(k+1)\)} & M1 & {States the assumption that \(k^5-k\) is divisible by 5 and considers the expression with \(n=k+1\).} \\
& {\((k+1)^5-(k+1)-(k^5-k)\)\\[4pt] \(\phantom{(}=(k^5+5k^4+\ldots+1)-(k+1)-(k^5-k)\)} & M1 & {Expands \((k+1)^5\) to six terms, condoning coefficient slips, and forms the difference, or another multiple of \(k^5-k\) subtracted, or rearranges to isolate \(k^5-k\).} \\
& {\(\phantom{(}=5k^4+10k^3+10k^2+5k\)} & A1 & {Or the correct equivalent expression for their route.} \\
& {\((k+1)^5-(k+1)\)\\[4pt] \(\phantom{(}=(k^5-k)+5(k^4+2k^3+2k^2+k)\)\\[4pt] \(\Rightarrow\) divisible by 5} & A1 & {cso. Needs the factor 5 shown and the assumption used for \(k^5-k\).} \\
& {True for \(n=1\); true for \(n=k\)\\[4pt] \(\Rightarrow\) true for \(n=k+1\)\\[4pt] \(\Rightarrow\) true for all positive integers \(n\)} & E1 & {States true for \(n=1\), true for \(n=k\) implies true for \(n=k+1\), hence true for all positive integers \(n\).} \\
\nopagebreak \SetCell[c=2]{r} Total for Question 1 & & 6 & \\
\end{msscheme}

\msQuestion{2}
\begin{msscheme}{
  hline{3-7} = {2-Z}{0.4pt, msGrid},
  hline{8} = {0.4pt, msGrid},
  row{8} = {font=\bfseries},
}
Question & Scheme & Marks & Notes \\
2 & {\(u_1 = 5^2 - 2^1 = 23\)} & B1 & {\(u_1\) evaluated as 23, or \(25 - 2\); ``true for \(n=1\)'' alone is B0.} \\
& {Assume \(5^{2k} - 2^k = 23m\)\\[4pt] \(u_{k+1} = 5^{2(k+1)} - 2^{k+1}\)} & M1 & {Assumes \(u_k\) divisible by 23 and writes \(u_{k+1}\), or \(u_{k+1} - \lambda u_k\), in terms of \(k\).} \\
& {\(\phantom{u_{k+1}} = 25 \times 5^{2k} - 2 \times 2^k\)} & M1 & {Uses index laws to write \(u_{k+1}\) in terms of \(5^{2k}\) and \(2^k\). Condone one coefficient slip.} \\
& {\(\phantom{u_{k+1}} = 23 \times 5^{2k} + 2\left(5^{2k} - 2^k\right)\)} & A1 & {oe, e.g. \(25\left(5^{2k}-2^k\right) + 23 \times 2^k\), or \(u_{k+1} - 2u_k = 23 \times 5^{2k}\).} \\
& {\(\phantom{u_{k+1}} = 23 \times 5^{2k} + 2(23m)\)\\[4pt] \(\phantom{u_{k+1}} = 23\left(5^{2k} + 2m\right)\)} & A1 & {cso. Uses the hypothesis to show \(u_{k+1}\) is a multiple of 23: factor 23 explicit, or each term stated divisible by 23.} \\
& {True for \(n = 1\); true for \(n = k\) \(\Rightarrow\) true for \(n = k+1\); so true for all \(n \in \mathbb{Z}^+\)} & E1 & {States all three points: true for \(n=1\), truth for \(n=k\) implies truth for \(n=k+1\), so true for all positive integers \(n\).} \\
\nopagebreak \SetCell[c=2]{r} Total for Question 2 & & 6 & \\
\end{msscheme}

\msQuestion{3}
\begin{msscheme}{
  hline{3-7} = {2-Z}{0.4pt, msGrid},
  hline{8} = {0.4pt, msGrid},
  row{8} = {font=\bfseries},
}
Question & Scheme & Marks & Notes \\
3 & {\(n=1:\ 9-8-1=0\), divisible by 64} & B1 & {Needs the value 0 and a statement that it is divisible by 64, or \(0=64\times0\).} \\
& {Assume \(9^k-8k-1=64m\)\\[4pt] Considers \(9^{k+1}-8(k+1)-1\)} & M1 & {States the assumption for \(n=k\) and writes the expression for \(n=k+1\). Not for assuming the result for \(n=k+1\).} \\
& {\(9^{k+1}-8(k+1)-1=9\times 9^k-8k-9\)\\[4pt] \(\phantom{9^{k+1}-8(k+1)-1}=9(9^k-8k-1)+\ldots\)} & M1 & {Uses \(9^{k+1}=9\times9^k\) and relates the result to \(9^k-8k-1\) or \(64m\), e.g. by substituting \(9^k=64m+8k+1\) or forming a difference with the \(n=k\) expression. Condone slips.} \\
& {\(\phantom{9^{k+1}-8(k+1)-1}=9(9^k-8k-1)+64k\)} & A1 & {Correct expression, oe such as \(9(64m+8k+1)-8k-9\).} \\
& {\(\phantom{9^{k+1}-8(k+1)-1}=9(64m)+64k\)\\[4pt] \(\phantom{9^{k+1}-8(k+1)-1}=64(9m+k)\)} & A1 & {cso. Shown as 64 times an integer expression, or both terms identified as multiples of 64.} \\
& {True for \(n=1\); true for \(n=k\Rightarrow\) true for \(n=k+1\)\\[4pt] so true for all positive integers \(n\)} & E1 & {Concludes, naming: true for \(n=1\); true for \(n=k\) implies true for \(n=k+1\); so true for all positive integers \(n\).} \\
\nopagebreak \SetCell[c=2]{r} Total for Question 3 & & 6 & \\
\end{msscheme}

\msQuestion{4}
\begin{msscheme}{
  hline{3-7} = {2-Z}{0.4pt, msGrid},
  hline{8} = {0.4pt, msGrid},
  row{8} = {font=\bfseries},
}
Question & Scheme & Marks & Notes \\
4 & {\(n=0\): \(2^0+2\times3^0+4^1=7\)} & B1 & {Evaluation to 7 shown. Checking \(n=1\) alone does not score.} \\
& {Assumes \(2^k+2\times3^{2k}+4^{2k+1}=7m\)\\[4pt] \(n=k+1\): \(2^{k+1}+2\times3^{2k+2}+4^{2k+3}\)} & M1 & {States the assumption for \(n=k\) and writes the expression with \(n=k+1\) substituted.} \\
& {\(= 2\times2^k+18\times3^{2k}+16\times4^{2k+1}\)} & M1 & {Writes each term of the \(n=k+1\) expression, or of a combination such as its difference from the \(n=k\) expression, as a multiple of \(2^k\), \(3^{2k}\) or \(4^{2k+1}\). Condone a coefficient slip.} \\
& {\(= 2\left(2^k+2\times3^{2k}+4^{2k+1}\right)\)\\[4pt] \(\phantom{=}+14\times3^{2k}+14\times4^{2k+1}\)} & A1 & {oe: a multiple of the \(n=k\) expression plus terms with a visible factor 7, or with one power replaced using the assumption.} \\
& {\(= 14m+14\left(3^{2k}+4^{2k+1}\right)\)\\[4pt] \(= 7\left(2m+2\times3^{2k}+2\times4^{2k+1}\right)\)} & A1 & {cso. The \(n=k+1\) expression as 7 times an integer expression, or a sum of multiples of 7, with the assumption applied.} \\
& {True for \(n=0\)\\[4pt] True for \(n=k\) \(\Rightarrow\) true for \(n=k+1\)\\[4pt] \(\Rightarrow\) true for all non-negative integers \(n\)} & E1 & {States true for \(n=0\), true for \(n=k\) implies true for \(n=k+1\), so true for all non-negative integers \(n\). Condone ``for all \(n\)''.} \\
\nopagebreak \SetCell[c=2]{r} Total for Question 4 & & 6 & \\
\end{msscheme}

\msQuestion{5}
\begin{msscheme}{
  hline{3-7} = {2-Z}{0.4pt, msGrid},
  hline{8} = {0.4pt, msGrid},
  row{8} = {font=\bfseries},
}
Question & Scheme & Marks & Notes \\
5 & {\(n=1: 7^1-6(1)-1=0\), divisible by 36} & B1 & {Needs the value 0 and a statement that it is divisible by 36, or \(0=36\times0\).} \\
& {Assume \(7^k-6k-1=36m\)\\[4pt] \(n=k+1: 7^{k+1}-6(k+1)-1\)} & M1 & {States the assumption for \(n=k\), in words or as \(=36m\), and writes the expression with \(n=k+1\).} \\
& {\(7^{k+1}-6(k+1)-1=7\times7^k-6k-7\)\\[4pt] \(\phantom{7^{k+1}-6(k+1)-1}=7(7^k-6k-1)+\ldots\)} & M1 & {Writes \(7^{k+1}\) as \(7\times7^k\) and expresses the \(n=k+1\) expression using \(7^k-6k-1\) or \(36m\), or subtracts a multiple of it. Condone slips in the remainder.} \\
& {\(\phantom{7^{k+1}-6(k+1)-1}=7(7^k-6k-1)+36k\)} & A1 & {Correct unsimplified expression, oe such as a correct difference equal to \(36k\).} \\
& {\(\phantom{7^{k+1}-6(k+1)-1}=7(36m)+36k=36(7m+k)\)} & A1 & {Factor of 36 shown, or every term stated to be a multiple of 36, using the assumption.} \\
& {True for \(n=1\); true for \(n=k\) \(\Rightarrow\) true for \(n=k+1\); so true for all \(n\in\mathbb{Z}^+\)} & E1 & {States true for \(n=1\), true for \(n=k\) implies true for \(n=k+1\), and hence true for all positive integers \(n\).} \\
\nopagebreak \SetCell[c=2]{r} Total for Question 5 & & 6 & \\
\end{msscheme}

\msQuestion{6}
\begin{msscheme}{
  hline{3-7} = {2-Z}{0.4pt, msGrid},
  hline{8} = {0.4pt, msGrid},
  row{8} = {font=\bfseries},
}
Question & Scheme & Marks & Notes \\
6 & {\(f(1) = 9 - 8 + 63 = 64\)} & B1 & {64 must be seen; \(f(1)\) stated divisible by 64 without its value is B0.} \\
& {Assume \(9^k - 8k + 63 = 64m\)\\[4pt] \(f(k+1) = 9^{k+1} - 8(k+1) + 63\)} & M1 & {States the assumption for \(n = k\) and substitutes \(k+1\) into \(f\). Not for assuming the result for \(n = k+1\).} \\
& {\(f(k+1) = 9 \times 9^k - 8k + 55\)\\[4pt] \(\phantom{f(k+1)} = 9(9^k - 8k + 63) + \ldots\)} & M1 & {Writes \(9^{k+1}\) as \(9 \times 9^k\) and expresses \(f(k+1)\) in terms of \(f(k)\), or forms \(f(k+1) - 9f(k)\) oe. Condone arithmetic slips.} \\
& {\(f(k+1) = 9(9^k - 8k + 63) + 64k - 512\)} & A1 & {oe, unsimplified accepted, such as \(f(k+1) - 9f(k) = 64k - 512\).} \\
& {\(f(k+1) = 9(64m) + 64(k-8)\)\\[4pt] \(\phantom{f(k+1)} = 64(9m + k - 8)\)} & A1 & {cso. Every term shown to be a multiple of 64, with \(f(k+1)\) stated divisible by 64.} \\
& {True for \(n = 1\);\\[4pt] true for \(n = k \Rightarrow\) true for \(n = k+1\);\\[4pt] so true for all \(n \geq 1\)} & E1 & {States all three points: true for \(n = 1\), true for \(k\) implies true for \(k+1\), so true for all \(n \geq 1\).} \\
\nopagebreak \SetCell[c=2]{r} Total for Question 6 & & 6 & \\
\end{msscheme}

\msQuestion{7}
\begin{msscheme}{
  hline{3-7} = {2-Z}{0.4pt, msGrid},
  hline{8} = {0.4pt, msGrid},
  row{8} = {font=\bfseries},
}
Question & Scheme & Marks & Notes \\
7 & {\(n=1\colon 7^2+8\times 13-9 = 144 = 36\times 4\)} & B1 & {Shows 144 and that it is a multiple of 36.} \\
& {Assumes \(7^{2k}+8\times 13^k-9 = 36m\)\\[4pt] \(n=k+1\colon 7^{2k+2}+8\times 13^{k+1}-9\)} & M1 & {States the assumption for \(n=k\) and writes the expression with \(n=k+1\). Condone divisibility assumed in words.} \\
& {\(\phantom{n} = 49\times 7^{2k}+104\times 13^k-9\)\\[4pt] \(\phantom{n} = 13(7^{2k}+8\times 13^k-9)+\ldots\)} & M1 & {Uses index laws and rearranges in terms of the \(n=k\) expression, or forms the \(n=k+1\) expression minus a multiple of it. Condone arithmetic slips.} \\
& {\(\phantom{n} = 13(7^{2k}+8\times 13^k-9)+36(7^{2k}+3)\)} & A1 & {Correct, unsimplified, oe such as \(49(7^{2k}+8\times 13^k-9)-36(8\times 13^k-12)\).} \\
& {\(\phantom{n} = 36(13m+7^{2k}+3)\)} & A1 & {Uses the assumption to show a multiple of 36, by a factor of 36 shown or each term stated divisible by 36.} \\
& {True for \(n=1\); true for \(n=k\)\\[4pt] \(\Rightarrow\) true for \(n=k+1\); so true for all\\[4pt] positive integers \(n\)} & E1 & {Concludes with all three: true for \(n=1\), true for \(n=k\) implies true for \(n=k+1\), and so true for all positive integers \(n\).} \\
\nopagebreak \SetCell[c=2]{r} Total for Question 7 & & 6 & \\
\end{msscheme}

\msQuestion{8}
\begin{msscheme}{
  hline{3-7} = {2-Z}{0.4pt, msGrid},
  hline{8} = {0.4pt, msGrid},
  row{8} = {font=\bfseries},
}
Question & Scheme & Marks & Notes \\
8 & {\(n=1: 13+12+11=36=12\times3\)} & B1 & {Needs 36 shown as a multiple of 12, or stated divisible by 12.} \\
& {Assumes \(13^k+3\times4^k+11=12m\)\\[4pt] Considers \(13^{k+1}+3\times4^{k+1}+11\)} & M1 & {States the assumption for \(n=k\) and writes the expression for \(n=k+1\). Not for the \(n=k+1\) expression alone.} \\
& {\((13^{k+1}+3\times4^{k+1}+11)\)\\[4pt] \(\phantom{=}-(13^k+3\times4^k+11)\)\\[4pt] \(=13^k(13-1)+3\times4^k(4-1)\)} & M1 & {Subtracts the \(n=k\) expression, or any multiple of it, from the \(n=k+1\) expression and writes \(13^{k+1}\), \(4^{k+1}\) in terms of \(13^k\), \(4^k\). Condone slips.} \\
& {\(=12\times13^k+9\times4^k\)} & A1 & {Or the correct simplified result for another multiple subtracted.} \\
& {\(13^{k+1}+3\times4^{k+1}+11\)\\[4pt] \(\phantom{1}=12m+12\times13^k+12(3\times4^{k-1})\)\\[4pt] \(\phantom{1}=12(m+13^k+3\times4^{k-1})\)} & A1 & {cso. Shows \(9\times4^k\) as a multiple of 12 and uses the hypothesis to write the \(n=k+1\) expression as 12 times an integer.} \\
& {True for \(n=1\); true for \(n=k\Rightarrow\) true for \(n=k+1\)\\[4pt] so true for all \(n\in\mathbb{Z}^+\)} & E1 & {Concludes with true for \(n=1\), true for \(n=k\) implies true for \(n=k+1\), and so true for all positive integers \(n\).} \\
\nopagebreak \SetCell[c=2]{r} Total for Question 8 & & 6 & \\
\end{msscheme}

\end{document}
